\section{Projekt Requirements}
\label{sec:Projekt Requirements}
\begin{longtable}{|p{1.7cm}|p{3.5cm}|p{7.2cm}|p{1.8cm}|p{2.2cm}|}
	\caption{Übersicht der Anforderungen für das Projekt}
	\label{tab:anforderungen} \\
	\hline
	\textbf{ID} & \textbf{Kategorie} & \textbf{Anforderungsbeschreibung} & \textbf{Priorität} & \textbf{Ersteller} \\
	\hline
	\endfirsthead
	
	\hline
	\textbf{ID} & \textbf{Kategorie} & \textbf{Anforderungsbeschreibung} & \textbf{Priorität} & \textbf{Ersteller} \\
	\hline
	\endhead
	
	\hline
	\multicolumn{5}{r}{Fortsetzung auf nächster Seite} \\
	\endfoot
	
	\hline
	\endlastfoot
	
	REQ-001 & Funktionale Anforderungen &
	Der Roboter soll sich vorwärts und rückwärts bewegen können. &
	High & Faik \\
	\hline
	REQ-002 & Funktionale Anforderungen &
	Der Roboter soll sich über eine Fernsteuerung (ESP32 + Fernbedienung oder App) steuern lassen. &
	High & Faik \\
	\hline
	REQ-003 & Funktionale Anforderungen &
	Der Roboter soll gleichmäßig links und rechts drehen können. &
	Medium & Faik \\
	\hline
	REQ-004 & Funktionale Anforderungen &
	Das Display soll Statusinformationen wie Akkustand und Verbindung anzeigen können. &
	Medium & Simon \\
	\hline
	REQ-005 & Nicht-funktionale Anforderungen &
	Die Motorsteuerung soll eine Verzögerung von weniger als 0,5 Sekunden haben. &
	High & Faik \\
	\hline
	REQ-006 & Nicht-funktionale Anforderungen &
	Die Motoren sollen mindestens 7\,kg Nutzlast sicher bewegen können. &
	High & Faik \\
	\hline
	REQ-007 & Nicht-funktionale Anforderungen &
	Die Geschwindigkeit des Roboters soll max.\ 10\,km/h sein. &
	Medium & Faik \\
	\hline
	REQ-008 & Nicht-funktionale Anforderungen &
	Die Motorsteuerung soll bei mindestens 10\,m Entfernung funktionieren können. &
	Low & Faik \\
	\hline
	REQ-009 & Safety &
	Wird kein Signal fürs Fahren erteilt, stoppt das Fahrzeug. &
	High & Angus \\
	\hline
	REQ-010 & Design &
	Display soll Augen anzeigen. &
	Low & Angus \\
	\hline
	REQ-011 &  &
	Kamera soll Bilddaten für Mikrocomputer liefern. &
	Medium & Angus \\
	\hline
	REQ-012 & Qualität &
	Motor soll genug Drehmoment haben, um das Fahrzeug fortzubewegen. &
	High & Angus \\
	\hline
	REQ-013 & Funktionale Anforderungen &
	Der Roboter soll in Fahrtrichtung lenken und Kurven fahren können. &
	High & Hai Nam \\
	\hline
	REQ-014 & Funktionale Anforderungen &
	Der Roboter soll sich um die eigene Achse drehen können. &
	High & Hai Nam \\
	\hline
	REQ-015 & Funktionale Anforderungen &
	Der Roboter muss sanft beschleunigen und verlangsamen. &
	Medium & Hai Nam \\
	\hline
	REQ-016 & Safety &
	Geschwindigkeit muss für Personen ungefährlich sein (4–5\,km/h). &
	High & Angus \\
	\hline
	REQ-017 & Safety &
	Der Roboter muss bei Hindernissen automatisch anhalten. &
	Low & Hai Nam \\
	\hline
	REQ-018 & Qualität &
	Die Struktur muss das gesamte Gewicht sicher tragen können. &
	High & Angus \\
	\hline
	REQ-019 & Design &
	Der Kopf soll sich horizontal drehen können. &
	High & Angus \\
	\hline
	REQ-020 & Qualität &
	Die Elektronik darf bei hoher Belastung nicht durchbrennen. &
	High & Angus \\
	\hline
	REQ-021 & Qualität &
	Die tragende Struktur soll Kräfte- und Geschwindigkeitsänderungen aushalten. &
	High & Angus \\
	\hline
	REQ-022 &  &
	Der Kopf soll sich vertikal drehen können. &
	Low & Angus \\
	\hline
	REQ-023 & Wartbarkeit &
	Einfache/selbsterklärende Bedienung. &
	Medium & Angus \\
	\hline
	REQ-024 & Wartbarkeit &
	Akku soll sich sicher wechseln lassen. &
	High & Angus \\
	\hline
	REQ-025 & Wartbarkeit &
	Elektronik darf beim Akkuwechseln nicht beschädigt werden. &
	Medium & Angus \\
	\hline
	REQ-026 & Safety &
	Innere Elektronik soll vor unbefugtem Zugriff geschützt sein. &
	Medium & Angus \\
	\hline
	REQ-027 & Wartbarkeit &
	Innere Elektronik soll vor Dreck und Umwelt geschützt sein. &
	Medium & Angus \\
	\hline
	REQ-028 & Safety &
	Erschütterungen sollen Elektronik nicht beschädigen. &
	Medium & Angus \\
	\hline
	REQ-029 & Wartbarkeit &
	Der Antrieb darf keinen hohen Verschleiß aufweisen. &
	Medium & Angus \\
	\hline
	REQ-030 & Design &
	Das Erscheinungsbild soll der DHBW-Designsprache entsprechen. &
	Low & Angus \\
	\hline
	REQ-031 & Design &
	Das Display soll animierte Augenbewegungen darstellen. &
	Low & Simon \\
	\hline
	REQ-032 & Funktionale Anforderungen &
	Der Kopf soll sich mindestens 180° drehen können. &
	Medium & Simon \\
	\hline
	REQ-033 & Nicht-funktionale Anforderungen &
	Servomotoren sollen eine Positionsgenauigkeit von ±2° erreichen. &
	High & Simon \\
	\hline
	REQ-034 & Funktionale Anforderungen &
	Raspberry Pi soll Display, Audio und KI-Funktionen ausführen. &
	Medium & Simon \\
	\hline
\end{longtable}
